\renewcommand{\chaptertitle}{Interpreter Runtime}

%  Make first page footer:
\fancypagestyle{firststyle}{%
\fancyhf{}% Clear header/footer
\fancyhead{}
\fancyfoot[L]{{\footnotesize
              %% We toggle between these:
              % Manuscript submitted for review.\\
              \authorname{}~\authorpatp{}.  ``\chaptertitle{}''.  {\it The Nock Book} (2025):  10–$x$.
              }}
}
%  Arrange subsequent pages:
\fancyhf{}
\fancyhead[LE]{{\urbitfont The Nock Book}}
\fancyhead[RO]{\chaptertitle{}}
\fancyfoot[LE,RO]{\thepage}

\chapter{\chaptertitle}

\thispagestyle{firststyle}

\begin{abstract}

\end{abstract}

Now that we have established a firm foundation for idiomatic programming in Nock, we can build an interpreter using a higher-level language to execute Nock formulas.  At this point, we will build the simplest possible Nock interpreter, a tree-walking interpreter.  This requires a cursor and some basic memory management but not much more.  We will implement some minimalist static hints to permit side effects like output.
